\chapter{Conclusions and Future Work}
\label{ch:conclusions}

\Cref{sec:problem-statement} asked under which conditions a neurosymbolic detector is worth its
extra complexity. On the three datasets and the knowledge bases of this thesis, the answer is
narrow. Domain knowledge used as a constraint in training gave no demonstrated gain in detection.
The same knowledge read at inference mattered in one case: a rule that states an attack mechanism,
on a dataset the models had never seen.

% ================================================================
% 11.1 Summary of the work
% ================================================================
\section{Summary of the work}
\label{sec:conc-summary}

The thesis built a benchmark for machine-learning and neurosymbolic detectors and ran it on NSL-KDD,
CIC-IDS2017 and CIC-DDoS2019. LTN-full, a network trained with domain axioms, faces the same network
trained without them, LTN-sup, and with cross-entropy, MLP-CE with and without class weights. Random
Forest, XGBoost, the axioms alone and a 1D-CNN complete the arms. Each arm with hyperparameters had
20 tuning trials, the learning arms ran on ten seeds, and five paired comparisons per dataset form
the confirmatory family. Stress experiments changed the number of labels, the rare classes, the attacks seen
in training, the dataset and the knowledge.

Of the contributions of \cref{sec:contributions}, C1 is delivered for the batch setting, and its
usability axis rests on an estimated effort and a scorecard. C2 and C3 are delivered, and most
of the results of C3 are negative. C4 has no measurements yet. C5, axiom-level satisfaction and the
domain-rule violation rate, was measured for every arm: a low violation rate proved not to be a merit, and the rule precision read beside it was the informative
figure.

% ================================================================
% 11.2 Answers to the research questions
% ================================================================
\section{Answers to the research questions}
\label{sec:conc-answers}

\paragraph{RQ1, detection quality.}
Not in these runs. With architecture, features and tuning budget matched, the domain knowledge has
no demonstrated effect on the macro-F1. LTN-full minus LTN-sup is $+0.008$, $+0.015$ and $+0.002$ on
NSL-KDD, CIC-IDS2017 and CIC-DDoS2019, positive on all three and significant on none after Holm
(\cref{sec:res-significance}). No class shows a gain. The LTN training without axioms has a cost of
its own: LTN-sup is below MLP-CE with weights by 0.0473 on NSL-KDD and 0.0206 on CIC-IDS2017. The
architecture does no better, since the 1D-CNN is below the MLP on NSL-KDD and equal on the CIC
datasets. On CIC-IDS2017 both trees stay ahead of every network.

\paragraph{RQ2, operating regime.}
The advantage does not appear under the conditions tested, with one partial exception. With few
labels the knowledge never helps: where the effect stands out of the seeds it is negative and at
most 1.5 points, so H4 is not supported (\cref{sec:res-scarcity}). With thinned rare classes it
gains nothing systematic, so neither is H2, and on attacks held out of training it changes nothing.
A larger weight on the axioms does not help: on NSL-KDD it lowers the macro-F1 by 0.015 to 0.023.
On CIC-IDS2017 every other weight stays within 0.004. Shifted thresholds and a swapped target change LTN-full by at most
about a point, on a grid that does not test H6 as written. Under noise LTN-full and LTN-sup lose
the same. What moves the macro-F1 is the exponent of the loss, for LTN-sup as for LTN-full.

The exception is the transfer from CIC-IDS2017 to CIC-DDoS2019 (\cref{sec:res-unseen}). The learned
models detect 0 to 13\,\% of the attacks, with a best mean macro-F1 of 0.4812, and the rules extracted
from the source never help. The mechanism axiom of the UDP flood lifts LTN-full by about ten points
in training, only to the level of MLP-CE with weights. Read alone at inference it reaches 0.918, and
no model with the post-check goes beyond it. The knowledge that transfers is that of a mechanism,
mostly at inference, for this pair of datasets, this direction and this axiom.

\paragraph{RQ3, cost.}
Partly answered. Each axiom added to the base of LTN-full costs about 2 to 4\,\% of the time of a
training epoch, with a growth close to linear. At inference an LTN-trained network costs as an MLP,
1.3 to 2.2\,ms per call with one row on the CPU of the laptop L1, and the post-check adds up to
2.2\,ms, on the CIC datasets less than the spread of the measurement (\cref{sec:res-cost}). Both halves of H7 agree with these runs. The total training cost of
the arms is not compared, since the final runs ran as parallel processes. The streaming half of the question
has no result yet.
% FIXME(open): E14 and E15 have no results; when they arrive, answer the real-time half of RQ3 here and update C4 in Sec. 11.1, the last item of Sec. 11.4 and Sec. 11.5.6.

\paragraph{RQ4, usability.}
One axiom is estimated at hours of work, a knowledge base at one or two days and the change of a rule
at an hour at most. The dearest part is the search for axioms that do not exist, which the
baselines do not pay. In return, the explanations cover one alert in three on NSL-KDD and eight to
nine in ten on the CIC datasets, where they never contradict the prediction
(\cref{sec:res-usability}). In enforce mode with the main bases the post-check lowers every learned
arm on the three datasets. It helps only where the rule is more precise than the model on the rows
where it fires, as in the transfer. On the scorecard LTN-full shares the lowest setup
score with the rule-only classifier and is the only arm that needs a new training run to change a
rule. It explains no better than the other learned arms, since the explanation comes from the post-check that every arm shares (\cref{tab:res-scorecard}).

\paragraph{RQ5, methodology.}
A fair benchmark here needs one network for the neural arms, equal tuning trials, selection on
validation data only, repeated seeds and a family of tests fixed before the runs
(\cref{sec:fair-benchmark}). Threshold-free measures, rule precision beside the violation rate,
counts of the overrides that correct and break a prediction, and per-seed results each changed the
reading of a result. The protocol also showed a limit of its own: on NSL-KDD the selection on a cut
of the training file does not carry to the test file, and this is reported, not corrected on the
test file. A threshold fixed on validation shows the same limit to an operator: the alert budget
holds on the test part of CIC-IDS2017 and breaks on the other two datasets
(\cref{tab:res-operating}).

% ================================================================
% 11.3 Practical guidelines
% ================================================================
\section{A decision checklist}
\label{sec:conc-checklist}

The questions below are for an engineer who holds a detection task and some rules. They come only
from the results above and share their limits.

\begin{enumerate}
  \item \textbf{Is a baseline already enough?} Train a Random Forest and a plain network first. On
    CIC-IDS2017 the trees beat every network, and on the CIC datasets no rule beat the models on its
    rows. If a baseline meets the need, stop here.
  \item \textbf{Is the rule a mechanism, and can it be read in your features?} An extracted rule
    fires almost only where the model is already right, as on CIC-IDS2017, and on the two datasets
    with a separate test part such rules fell to 0.130 and 0.547. A training confidence of 0.90 is no
    guarantee: A3 had it. Prefer a mechanism, and measure its precision on labelled data first: on
    single flows of CIC-IDS2017 a slow DoS axiom reached 0.008. Drop a rule that fails.
  \item \textbf{Is the rule more precise than the model where it fires?} Compare the two on those
    rows and count the overrides that correct and break a prediction. If the rule wins, it may change
    decisions. If not, use it to annotate alerts.
  \item \textbf{Do you need explanations, or transfer to new data?} The post-check gives
    explanations to any model, with no LTN training. A mechanism rule is worth most on new data, as
    in the transfer. Keep such a rule at inference.
  \item \textbf{Is a constraint in training justified?} On these runs, not for detection. Its cost
    per axiom is small, but LTN training without axioms lost to MLP-CE with
    weights on two datasets, and the one gain of the axioms only reached the best arm without
    knowledge. If you train with LTN, tune the exponent $p$.
  \item \textbf{Is the work justified?} Count hours per axiom, a day or two per base and the search
    for rules that may not exist. An attacker who knows a threshold evades the rule, so a decision
    that rests on a rule alone is exposed.
\end{enumerate}

On these results we would start from a tree or a plain network, add mechanism rules at inference
where they beat the model, and not train with logic for detection quality alone.

% ================================================================
% 11.4 Limitations
% ================================================================
\section{Limitations}
\label{sec:conc-limitations}

The limits of \cref{sec:threats,sec:res-threats} bound every statement of this chapter.

\begin{itemize}
  \item \emph{Knowledge.} The base of CIC-IDS2017 is extracted from the data in full, so there the
    effect of the knowledge is that of rules distilled from the labels the network also sees. The
    thresholds of the mechanism axioms were calibrated on training labels, except that of the UDP
    axiom of CIC-DDoS2019, which means no backward packet. That axiom entered the transfer from the
    training day of the target dataset. On NSL-KDD the axioms cover DoS and Probe only, and no usable
    axiom was found for R2L and U2R.
  \item \emph{Features.} On the CIC datasets every feature describes one flow, with no window or
    source statistic, so most candidate axioms of \cref{ch:attacks-ddos} could not be written there.
  \item \emph{Datasets.} NSL-KDD is old and synthetic, and its test file does not follow its
    training file. The test part of CIC-IDS2017 is in-distribution, CIC-DDoS2019 is close to
    saturation, and the transfer runs in one direction.
  \item \emph{Method.} Only LTN is tested. The other families of \cref{ch:learning-paradigms} are not
    run on traffic, and the maturity of their tools was not compared. LTN-sup and LTN-full were tuned
    separately, so their difference holds the knowledge together with its tuning.
  \item \emph{Measures.} The effort of E16 is an estimate. The times of E12 and E13 come from
    L1. Ten seeds speak of the randomness of training on one dataset. E14 and E15 have no results.
\end{itemize}

% ================================================================
% 11.5 Future work
% ================================================================
\section{Future work}
\label{sec:future-work}

The results leave six directions open. They concern richer axioms, learning the rules,
other frameworks, knowledge from threat intelligence, adversarial robustness and deployment
at line rate.

\subsection{Temporal and relational axioms}
\label{sec:fw-logics}

Mechanism axioms did not hold on single flows of CIC-IDS2017: the slow DoS axiom reached a precision
of 0.008 and the scan axiom 0.092 (\cref{sec:res-knowledge}). A scan or a flood is a pattern over
many flows, so the next step is window statistics, with one definition for training and streaming,
and axioms over sequences of flows and relations between sources and destinations.

\subsection{Learning the rules}
\label{sec:fw-learning-rules}

The rules that failed outside the training data are extracted ones, two of the four of this thesis.
A3 of NSL-KDD entered the base with a training confidence of 0.900 and fell to 0.130 on the test
file. Inductive logic programming or a better extractor would need another criterion. We would keep
a learned rule only if it holds on data its search did not see, from another day or dataset, and
would not learn it from the labels that train the detector.

\subsection{Other NeSy frameworks under the same benchmark}
\label{sec:fw-frameworks}

The negative result on training may belong to LTN, to its loss or to logic as a loss in general, and
one framework cannot separate them. Semantic Loss and DL2 would change the loss and keep the setting.
DeepProbLog and NeurASP place the network inside a logic program. In our reading the logic then
stays in the decision, where the knowledge acted in these runs, and they are the more interesting
arms. A new arm needs only the interface of a model
(\cref{sec:impl-models}) and inherits the splits, the tuning budget and the tests.

\subsection{Knowledge from threat intelligence}
\label{sec:fw-attack}

The search for axioms by hand is the dearest part of the knowledge, and \cref{sec:alert-formats}
names ATT\&CK and threat-intelligence feeds as a possible source. As we read them, they describe the
behaviour of attacks and not thresholds on columns, so the open step is the map from a technique to a condition on
flow features. Such knowledge exists before the dataset, unlike the base of CIC-IDS2017, and the
transfer is ready as its test.

\subsection{Adversarial robustness}
\label{sec:fw-adversarial}

E10 assumes an adversary who knows the thresholds. It broke the rule-only classifier on the three
datasets, while no network reacted (\cref{sec:res-noise}). The
knowledge that helps at inference is therefore the part exposed to evasion. E10 does not ask whether
a moved flow still works as an attack. A study of an adaptive adversary, left out by
\cref{sec:scope-limitations}, would limit the perturbation to what keeps the attack alive, test model
and post-check together, and compare mechanism rules with extracted ones.

\subsection{Deployment at line rate}
\label{sec:fw-line-rate}

The times of E12 are calls on the CPU of L1 and say nothing on line rate. An LTN-trained
network is a plain network at inference, so quantization applies as to any MLP, and the rules
could move to an eBPF or DPDK front end. The
features come from closed flows, however, so such a front end needs rules on packets or on partial
flows. The end-to-end measurements of E14 and E15 have to come first.
